\section{Analysis Tab}
\label{sec:analysis}

The \app{Analysis} tab is where measured data are fitted to a
hyperfine model. Each fit lives in its own \emph{analysis project}
sub-tab, and a permanent \app{Isotope Shifts} sub-tab combines the
fitted centroids across projects into shift tables and plots.

\subsection{Projects and the block pipeline}

\app{+ Create Analysis Project} is a menu with two entries:
\app{Sample Project} (the standard fit pipeline; the isotopes you
actually want to measure) and \app{Reference Project} (identical
pipeline, but its fitted centroids feed the Gaussian-process drift
correction of Section~\ref{sec:analysis_gp} instead of the shift
table). Right-click a project tab to convert between the two kinds,
rename, or close it. A project chains four kinds of block, drawn
left-to-right with a colour-coded border:

\begin{description}
    \item[\app{Source} (blue).] Loads the run files, applies gates,
        and bins to voltage or rest-frame frequency.
    \item[\app{Model} (green).] Defines a satlas2 model
        (one HFS per isotope or isomer, plus optional sidepeaks and
        background/auxiliary models).
    \item[\app{Fitter} (orange).] Picks the optimiser, statistic,
        sharing, constraints, and runs the fit.
    \item[\app{Output} (purple).] Selects which reports, plots, and
        diagnostics are produced and pushed to the \app{Results} tab.
\end{description}

\app{+ Add Block} inserts a new block of any type; the hamburger
handle on each block re-orders them by drag, the title checkbox
disables a block, and dragging a block's right edge resizes it.

\subsection{Source block}

The \app{Files} list at the top of the source block holds the run
files attached to this source. \app{+ Add Files} opens an
ASDF/\code{.vasdf} file picker; \app{Import from Pre-Analysis} copies
the ticked files, physics parameters, gates, and cooler/laser
overrides from a pre-analysis project in one step;
\app{Merge Selected} produces a synthetic merged entry analogous to
the merge in the \app{Pre-Analysis} tab.

\app{Physics Parameters} ($Z$, $A$, mass, harmonic, lower / upper
levels) feed the per-run rest-frame conversion. The mass field shows
the AME-table value; tick \app{Override} to enter a custom value.

\app{Gates} are the standard set: TOF window in microseconds, PMT
channel mask plus an optional DC channel, lab-frame voltage gate,
rest-frame frequency gate (always entered in MHz), and a noise-filter
threshold (maximum events per timestamp; 0 disables).

\paragraph{Binning.}
\app{Bin mode} chooses the spectrum domain --- \app{Frequency} (rest
frame) or \app{Raw Voltage} (lab frame, one bin per unique DAC
value). In frequency mode, \app{Bin definition} selects how the axis
is cut: \app{Per scan step} (the default --- one bin per native
voltage step, Doppler-projected; structurally immune to
uniform-grid aliasing), \app{Auto} (clstools' uniform grid at
roughly one bin per step --- can alias against the step spacing,
producing doubled-count spikes), \app{Fixed bin count}, or
\app{Fixed bin width [MHz]}. \app{Bin multiple} groups $N$ adjacent
native steps into one bin (per-step and raw-voltage binning). The
y-error model (\app{Poisson $\sqrt{y+1}$} default, with an exact
Poisson interval substituted below 10 counts), the x-column, and an
optional x-error from the per-step voltage spread complete the
policy.

Right-clicking a run row offers per-file \emph{binning overrides}
(freeze different binning for one run only, copy an override to the
checked runs, show the effective config), the binning diagnostics,
\app{Filter scans\dots}, and \app{Calibration\dots} --- the same
scan-filter and calibration machinery as Pre-Analysis
(Sections~\ref{sec:pa_scanfilter}, \ref{sec:pa_calibration}), backed
by the same shared registries.

\paragraph{Binning dialog.}
The \app{Binning} button opens a three-tab health check of the actual
binning for every checked file. \app{Summary} is a per-run table with
occupancy columns (empty bins with a scan-gap-aware denominator,
median and total counts). \app{Run detail} shows one run's spectrum
with its bin edges, an optional raw-event rug, a warning strip
(fallbacks, width mismatches, empty-bin and aliasing flags), and a
scale strip reporting the raw-voltage bin size, the measured
1~V~$\approx$~$N$~MHz conversion from the run's own events, and the
equivalent rest-frame bin width. \app{Compare} re-bins all runs onto
one grid and shows them as a heatmap.

\subsection{Model block}

A \app{Model} block builds one satlas2 model. \app{Type} offers
\app{HFS} (default), \app{Voigt}, \app{Skewed Voigt},
\app{Exponential Decay}, \app{Piecewise Constant}, and
\app{Polynomial} --- combine an HFS model with a Polynomial on the
same source for a sloping background. To seed parameters from the
trial overlay used in pre-analysis, click
\app{Import from Pre-Analysis}; \app{Retrieve Last Fit} re-imports
the most recent fit values to refine on top of them.

\paragraph{HFS settings.}
The HFS lineshape is a Voigt profile (the satlas2 version in use
supports only Voigt for HFS; use the dedicated \app{Skewed Voigt}
model type for asymmetric single peaks). Tick \app{Racah} (default)
to constrain peak amplitudes to their angular-momentum ratios;
untick it and use the per-component editor under
\app{Show Peak Amplitudes} when peaks deviate from Racah ratios ---
e.g.\ optical-pumping effects. \app{Sidepeaks} adds $N$
charge-exchange satellites at a fixed \app{Offset}, optionally with
Poisson-distributed amplitudes.

\paragraph{Parameters.}
The parameter table has five columns: \app{Parameter}, \app{Value},
\app{Bounds} (min, max), \app{Mode}, and \app{Expression}. There is
no separate vary checkbox --- the \app{Mode} combo is the single
source of truth:

\begin{description}
    \item[\app{Free}.] Floats during the fit (the only varying mode).
    \item[\app{Fixed}.] Held at its value. ($I$, $J_l$, $J_u$ are
        always fixed --- they are physics, not fit parameters.)
    \item[\app{Equal} / \app{Ratio} / \app{Offset}.] Tie the
        parameter to another via an auto-generated expression.
    \item[\app{Custom}.] A free-form expression typed in the
        \app{Expression} cell, validated live (red tint on errors)
        and again in a blocking pre-fit check that catches unknown
        names, self-references, and dependency cycles.
\end{description}

Model-block expressions use bare local names (\code{Al},
\code{centroid}, \dots) and are lowered to one full
\code{Run\_<id>\_\_\_<Model>\_\_\_<param>} name per source at fit
time; the Fitter block's \app{Advanced Constraints} use the full
names directly.

\subsection{Fitter block}

\app{Fit Mode} chooses how sources combine:

\begin{description}
    \item[\app{Separate (per file)}] (default). One independent fit
        per run, in parallel across worker processes; each run gets
        its own parameter copy.
    \item[\app{Simultaneous}.] One satlas2 fitter over all runs, and
        the \app{Parameter Sharing}, \app{Advanced Constraints}, and
        \app{Common Grid} groups appear. Sharing is \emph{opt-in}:
        tick a bare parameter and choose \app{Model} (share across
        same-named models) or \app{All} (share across every source).
        A \app{Common Grid} option re-bins all sources onto one
        shared x grid (width user-set or auto).
\end{description}

\app{Method} selects the optimiser (\code{leastsq},
\code{least\_squares}, \code{slsqp}, \code{cobyla}, \code{nelder},
\code{powell}, or the MCMC sampler \code{emcee}). \app{Statistics}
selects the objective: \app{Chi-square}, \app{Gaussian LLH}, or
\app{Poisson LLH} (recommended at low counts). \app{Scale
covariance} (off by default) multiplies parameter errors by
$\sqrt{\chi^2_{\mathrm{red}}}$; it is forced off for likelihood
fits.

\app{Show Priors} reveals a Gaussian-prior table (full parameter
names; each prior adds a penalty term) and \app{MCMC Settings}
exposes walkers, steps, burn-in, thinning, and auto-convergence ---
both take effect in \app{Simultaneous} mode (priors and advanced
constraints are ignored, with a pre-fit warning, in Separate mode).

\app{Run Fit} executes; the progress bar blends per-run completion
with in-run iteration progress. \app{Stop} cancels cooperatively;
\app{Revert} restores the model parameters to their pre-fit
snapshot. \app{Find Parameters} opens the Auto-Fitter
(Section~\ref{sec:autofitter}).

\subsection{Output block}

\app{Reports} produces the text fit report (with the correlation
matrix above a \app{min} threshold), plus \app{Parameters CSV} and
\app{Metadata CSV}. \app{Fit Plots} controls per-run figures:
\app{Individual fit plots}, optional \app{ToF plots} (ToF only /
ToF + spectrum / ToF + spectrum + fit), the \app{Residual panel},
\app{Hide zero-count bins}, \app{Component overlay} with
per-component checkboxes, and \app{Fit values on plot} --- an
on-plot values box (``Name = value $\pm$ err'') fed by a
per-parameter picker tree, re-rendered verbatim by the Results tab.
The \app{Plot Options\dots} button opens the per-plot-type styling
dialog (global rcParams, fit-plot shade-under-fit and values-box
size, tracker/walk/correlation/$\chi^2$-map styles, and the
isotope-shift plot).

\app{Tracker Plots} (Separate mode) plots chosen parameters against
\app{Run number} or \app{Start time}. \app{Diagnostics} adds the
$\chi^2$ map (any method) and, for \code{emcee}, correlation / walk
plots and confidence bands. \app{Iteration} numbers the output
folder (\app{Auto} increments \code{iter\_001},~\dots; \app{Manual}
uses a custom label).

\paragraph{Update iteration.}
The \app{Update iteration} button regenerates the \emph{last} fit's
iteration folder with the current output toggles, \emph{without}
re-running the fit --- useful when you realise a plot was left
unticked. Newly-ticked walk/correlation plots are rebuilt from the
saved MCMC chains; only the $\chi^2$ map and confidence bands need a
true re-run (the dialog says so rather than failing silently).

Results are written to
\code{<output\_root>/analysis/<project>/iter\_NNN/} --- the config
snapshot, report, CSVs (including a per-run binning summary and any
binning warnings), \code{plots/} with \code{.npz} sidecars for live
re-rendering, \code{tracking/}, \code{diagnostics/}, and
\code{chains/} for MCMC --- and forwarded to the \app{Results} tab.

\subsection{Auto-Fitter}
\label{sec:autofitter}

\app{Find Parameters} on the \app{Fitter} block opens the
Auto-Fitter. It runs a multi-start parallel \code{leastsq} sweep ---
\app{Starts} Gaussian perturbations of the current parameters (width
set by \app{Spread}; the first start is always the unperturbed
guess), each fitted in a worker process --- and ranks the results by
$\chi^2$ in a live table with a fit preview. \app{Refine} re-seeds a
new round around the selected result with the spread halved;
\app{Accept Selected} writes the chosen parameters back into the
Model blocks as one undoable step. The Auto-Fitter only seeds
initial guesses --- run the real fit afterwards. Fixed and
expression-tied parameters are never perturbed, and in simultaneous
mode the sweep respects the Fitter block's parameter sharing.

\subsection{Isotope Shifts sub-tab}

The permanent \app{Isotope Shifts} sub-tab combines fitted centroids
across sample projects into shift tables and a comparison plot. The
left panel holds a \app{Setup} tab and the
\app{Reference Correction (GP)} tab
(Section~\ref{sec:analysis_gp}); the right panel holds \app{Plot},
\app{Table}, and \app{Log}.

\app{Refresh from Projects} populates one row per sample project
with its label, mass number, run selection (\app{Weighted Average}
--- inverse-variance with Birge inflation when runs scatter --- or a
single run), centroid, and statistical error; the \app{Ref} radio
marks the reference isotope. A collapsible \app{Run preview} panel
under the table shows the highlighted run's spectrum, fit curve, and
reduced $\chi^2$. \app{Systematic Error} is either one manual MHz
value or an auto-estimate from the run-to-run jitter of cooler
voltage and laser setpoint (via the Doppler sensitivity).

\app{Compute Shifts} produces the plot and a 14-column table that
splits every uncertainty into its budget: per-isotope
$\sigma_{\mathrm{fit}}$ / $\sigma_{\mathrm{scatter}}$ /
$\sigma_{\mathrm{corr}}$ (from the GP) / $\sigma_{\mathrm{volt}}$,
and per-shift statistical, scatter, correction, systematic, and
total components. When two isotopes share the same GP corrector,
the correlated part of the correction cancels in the shift and the
cross-covariance is subtracted. $\delta\nu$ arrows are labelled in
the compact \code{value(err)~MHz} form. \app{Plot Options\dots}
configures layout (stacked / overlaid), content, x-axis, centroid
and reference lines, arrows, boxed labels, and colours.
\app{Send to Results} writes the study (CSV + log + live plot) to an
\code{IS\_<study>} iteration and forwards it to the \app{Results}
tab; \app{Export Table CSV} writes the table alone.

\subsection{Reference correction (Gaussian process)}
\label{sec:analysis_gp}

Slow drifts of the wavemeter or laser lock move every centroid
measured over a campaign. \denis{} implements the time-dependent
reference-centroid correction of van den Borne
(2025)~\cite{vandenborne2025}: fit the reference isotope's scans in
one or more \app{Reference Projects}, then, on the
\app{Reference Correction (GP)} tab, pick a kernel
(\app{RBF~+~WhiteNoise}, \app{Matérn(5/2)~+~WhiteNoise}, or the
composite thesis kernel with a periodic day/night term) and click
\app{Fit GP}. Each fitted reference centroid becomes one training
observation $(t, \mathrm{centroid}, \sigma)$; the diagnostic plot
shows the scatter, the MAP curve, and $1\sigma$/$2\sigma$ bands
(optionally with an MCMC posterior for diagnostics --- prediction
always uses the MAP point; PyMC~v5 is required).

\app{Compute per-file corrections} then fills a table with one row
per sample run: the correction $\mu_{\mathrm{GP}}(t)$, its
$\sigma$, and a $|\mu|/\sigma$ significance colouring. Each row's
\app{Mode} is \app{Auto} (use the GP), \app{Manual} (type a value),
or \app{Off}; the global \app{Apply correction at fit time} checkbox
gates the whole mechanism. Corrections are subtracted from the
binned MHz axis after the Doppler transform --- consistently in the
fit pipeline, the merge dialog, and the Auto-Fitter --- and the
propagated correction uncertainty (including the GP cross-covariance
between runs) feeds the $\sigma_{\mathrm{corr}}$ columns of the
shift table. The trained GP and the table round-trip through the
save file; corrections saved without their GP are dropped on load
with a prompt to re-fit.
